\chapter{Protoespacio sin lenguaje reciproco}
\label{ch:proto-sin-brillouin}

\section{El nudo conceptual}

El lenguaje de Brillouin es util como representacion espectral de modelos
periodicos homogeneos, pero no agota la pregunta ontologica por el
protoespacio. Identificar \Proto\ con el toro reciproco confunde un
\emph{dispositivo de calculo} con la \emph{estructura subyacente}.

\section{Direccion de trabajo}

\begin{enumerate}[label=(E\arabic*),leftmargin=*]
  \item Formular \Proto\ como tupla \ProtoTuple\ sobre grafos locales no
        necesariamente periodicos.
  \item Verificar que el sector efectivo Dirac sobrevive bajo
        perturbaciones que rompen la traslacion exacta pero preservan la
        localidad.
  \item Construir un test combinatorio (z3) sobre familias de grafos finitos
        donde se exija doblamiento fermionico de cargas quirales sin asumir
        toro reciproco.
\end{enumerate}
